\fchapter{Pág 96}

\fsection[\textcolor{waveBlue2}{Ejercicio 1}]{\boldit{\textcolor{waveBlue2}{TU SECTOR PROFESIONAL}}{
Piensa en un producto relacionado con las actividades de tu
sector profesional y explica las fases del modelo de economía lineal
por las que pasa (extracción de recursos para su fabricación,
producción, distribución, consumo y desecho).
}}

\begin{solution}

\boldit{Robot industrial de 6 ejes.} 

\begin{enumerate}
    \item \textbf{Extracción de recursos}
    \begin{itemize}
        \item Metales: acero, aluminio, cobre (estructura y cableado).
        \item Tierras raras para imanes de servomotores.
        \item Silicio y semiconductores para electrónica.
        \item Plásticos y polímeros para carcasas y aislamiento.
    \end{itemize}

    \item \textbf{Producción}
    \begin{itemize}
        \item Fabricación de piezas mecánicas (brazos, ejes, reductores).
        \item Ensamblaje de motores, sensores y controladores.
        \item Programación y pruebas.
        \item Integración en armarios eléctricos.
    \end{itemize}

    \item \textbf{Distribución}
    \begin{itemize}
        \item Transporte al integrador o cliente.
        \item Instalación en planta.
        \item Integración con PLC y SCADA.
    \end{itemize}

    \item \textbf{Consumo}
    \begin{itemize}
        \item Uso en tareas como soldadura o paletizado.
        \item Consumo eléctrico continuo.
        \item Mantenimiento preventivo y correctivo.
    \end{itemize}

    \item \textbf{Desecho}
    \begin{itemize}
        \item Fin de vida útil por desgaste u obsolescencia.
        \item Reciclaje de metales.
        \item Gestión de residuos electrónicos (RAEE).
        \item Parte no reciclable a vertedero.
    \end{itemize}
\end{enumerate}
\end{solution}
